\documentclass{article}
\usepackage[utf8]{inputenc}
\usepackage{amsmath}
\usepackage{amssymb}
\usepackage{booktabs}
\usepackage{geometry}
\geometry{a4paper, margin=1in}

\title{Maximal Independent Set (MIS) Algorithm Performance Report}
\author{MIS Project Team}
\date{\today}

\begin{document}

\maketitle

\section{Introduction}
This report evaluates various implementations of the Maximal Independent Set (MIS) algorithm. MIS is a fundamental problem in graph theory with applications in distributed computing, wireless networks, and resource allocation. We compare a sequential implementation against parallel versions based on Luby's algorithm.

\section{Algorithms}
The following algorithms were implemented and tested:
\begin{itemize}
    \item \textbf{Sequential}: The base $O(N+M)$ sequential implementation.
    \item \textbf{Luby}: A parallel implementation of Luby's algorithm with $O((N+M)/T \cdot \log(N))$ complexity.
    \item \textbf{Luby Improved}: An optimized version of Luby's algorithm utilizing dynamic OpenMP scheduling (\texttt{pragma omp parallel for schedule(dynamic, 32)}) to handle unbalanced node degrees.
\end{itemize}

\section{Experimental Setup}
Benchmarks were performed on two primary graph types:
\begin{enumerate}
    \item \textbf{Uniform Sparse}: Graphs with a low constant probability of edge existence.
    \item \textbf{Scale-Free}: Graphs generated using the Barabási–Albert model, resembling social network structures with power-law degree distributions.
\end{enumerate}
Each algorithm was tested using both standard Adjacency List (Graph) and Compressed Sparse Row (GraphCSR) formats.

\section{Results}
The following results were captured from the initial benchmark suite (5 graphs, 10 runs each, $N=1,000,000$).

\subsection{Standard Adjacency List Benchmarks}
\begin{center}
\begin{tabular}{l l r r r}
\toprule
Algorithm & Graph Type & Mean Time & Graph Std & Avg Run Std \\
\midrule
Sequential & Uniform Sparse & 8.837 ms & 0.368 ms & 0.417 ms \\
Luby 1t & Uniform Sparse & 209.911 ms & 0.721 ms & 3.407 ms \\
Luby 8t & Uniform Sparse & 46.948 ms & 1.792 ms & 1.675 ms \\
Luby Improved & Uniform Sparse & 27.178 ms & 2.910 ms & 1.929 ms \\
\midrule
Sequential & Scale-Free & 14.516 ms & 0.235 ms & 0.448 ms \\
Luby 1t & Scale-Free & 113.840 ms & 1.052 ms & 2.173 ms \\
Luby 8t & Scale-Free & 63.341 ms & 43.632 ms & 23.487 ms \\
Luby Improved & Scale-Free & 27.630 ms & 10.511 ms & 10.854 ms \\
\bottomrule
\end{tabular}
\end{center}

\subsection{CSR (Compressed Sparse Row) Benchmarks}
\begin{center}
\begin{tabular}{l l r r r}
\toprule
Algorithm & Graph Type & Mean Time & Graph Std & Avg Run Std \\
\midrule
Sequential CSR & Uniform Sparse\_csr & 9.806 ms & 1.121 ms & 0.305 ms \\
Luby 1t CSR & Uniform Sparse\_csr & 192.208 ms & 1.242 ms & 3.607 ms \\
Luby 8t CSR & Uniform Sparse\_csr & 42.141 ms & 0.312 ms & 1.634 ms \\
Luby Improved CSR & Uniform Sparse\_csr & 24.641 ms & 0.807 ms & 2.796 ms \\
\midrule
Sequential CSR & Scale-Free\_csr & 5.254 ms & 0.102 ms & 0.159 ms \\
Luby 1t CSR & Scale-Free\_csr & 73.781 ms & 1.570 ms & 2.005 ms \\
Luby 8t CSR & Scale-Free\_csr & 25.240 ms & 4.114 ms & 0.460 ms \\
Luby Improved CSR & Scale-Free\_csr & 14.763 ms & 1.609 ms & 2.322 ms \\
\bottomrule
\end{tabular}
\end{center}

\section{Conclusion}
The results indicate that while Luby's algorithm provides parallelization benefits, the Sequential implementation remains highly competitive for $N=1,000,000$ due to lower overhead. The "Improved" parallel version significantly outperforms the base Luby implementation by better managing workload distribution.

\section{Weighted MIS}

\subsection{Problem Formulation}

We recall that the standard Maximal Independent Set (MIS) is a subset of vertices $S\subset V$ of a graph $G=(V, E)$ such that no two vertices in $S$ are joined by an edge, and no new vertex can be added to $S$ without creating an edge. Now, we extend this framework by adding a weight $w_v\in \mathbb{R}^+$ for every vertex $v\in V$. In this setting, the goal is now to find a MIS with the maximum possible total weight \textbf{in practice}. It is important to note that finding the \textbf{maximum weight independent set} (i.e., the globally optimal solution) is NP-hard in general. In this view, we will focus on weight-based heuristic algorithms that find a valid MIS while favoring vertices with high weights, and then compare their solution quality empirically. 

\subsection{Algorithms}

We decided to study the performance of the following two algorithms, and then benchmark their performance based on the initial weight assignment. 

\subsubsection{Algorithm 1: Weighted Greedy}
A vertex $v$ is added to the MIS if and only if it has the highest weight among its active neighbors:
\begin{equation}
  v \text{ wins} \iff w_v \geq w_u \quad \forall\, u \in N(v),\, u \text{ active}
\end{equation}
Ties are broken by vertex index. This approach is based on Basagni~\cite{basagni2001}.
\subsubsection{Algorithm 2: Weighted Sampling}
Each active vertex $v$ draws $U_v \sim \text{Uniform}(0,1)$ and computes a priority:
\begin{equation}
  p_v = U_v^{1/w_v}
\end{equation}
Vertex $v$ wins the round if $p_v > p_u$ for all active neighbors $u$. That is, we have that the vertices with higher weights have the higher priority. 

\subsection{Initial Weight Distributions}

\textbf{Uniform:} Each vertex $v$ is assigned a weight drawn uniformly at random:
\begin{equation}
  w_v \sim \text{Uniform}(1, 10)
\end{equation}
\textbf{Exponential:} Each vertex $v$ is assigned:
\begin{equation}
  w_v = 1 + X, \quad X \sim \text{Exponential}(\lambda)
\end{equation}
The $+1$ shift ensures $w_v > 0$, which is required by the power trick in Algorithm 2.
\textbf{Clustered:} A fraction $\rho = 0.1$ of vertices are randomly designated as hotspots:
\begin{equation}
  w_v = \begin{cases} 10 & \text{with probability } \rho \\ 1 & \text{otherwise} \end{cases}
\end{equation}

\subsection{Experimental Results}

\subsection{Discussion}

\begin{thebibliography}{9}

    \bibitem{basagni2001}
    S. Basagni. Finding a maximal weighted independent set in wireless networks. Telecommunication Systems, 18(1–3):155–168, 2001. URL http://dx.doi.org/10.1023/A:1016747704458.
\end{thebibliography}

\end{document}
